% =====================================================================
%  MATH300 Project Description -- Continued Fractions
% =====================================================================
\documentclass[11pt,a4paper]{article}

\usepackage[T1]{fontenc}
\usepackage[utf8]{inputenc}
\usepackage{lmodern}
\usepackage[margin=2.5cm]{geometry}
\usepackage{enumitem}
\usepackage{booktabs}
\usepackage{tabularx}
\usepackage{array}
\usepackage{amsmath,amssymb}
\usepackage{xcolor}
\usepackage[colorlinks=true,urlcolor=blue,linkcolor=black,citecolor=black]{hyperref}

\setlist{itemsep=2pt, topsep=4pt}
\setlength{\parindent}{0pt}
\setlength{\parskip}{6pt}

% Placeholder macro: blue text in square brackets.
% To hide the colour in the final version, change {blue} to {black}.
\newcommand{\ph}[1]{\textcolor{blue}{[#1]}}

% Marker for optional bonus tasks.
\newcommand{\bonus}{\textbf{(Bonus)}}

\newcolumntype{L}{>{\raggedright\arraybackslash}X}

\begin{document}

% ---------------------------------------------------------------------
%  Title (no logo, no instructor name)
% ---------------------------------------------------------------------
\begin{center}
  {\large\bfseries MATH300: Mathematical Computing \textperiodcentered{} Fall 2026}\\[6pt]
  {\Large\bfseries Project: Continued Fractions for Constants and Functions}
\end{center}
\vspace{6pt}

Groups must be formed, and project titles submitted on Canvas, by
\textbf{October 26, 2026}. The project
has two graded deliverables and an oral exam: a 3--5 page report for
Project~2 (due \textbf{November 20, 2026}), a poster for the final project
(due \textbf{December 11, 2026}), and an oral exam held between November 30
and December 17, 2026, 3:30\,pm. Read the general rules below before you
start; they apply throughout the project.

% ---------------------------------------------------------------------
\section*{General remarks}
% ---------------------------------------------------------------------
\begin{itemize}
  \item Projects are set for a group of \textbf{two or three} students.
  \item Form your group by \textbf{October 26, 2026}, and submit your
        project title as a group assignment on \textbf{Canvas} by the same
        date.
  \item Every group member must carry an \textbf{equal portion} of the coding
        and of the report and poster preparation.
  \item The scope can be adjusted to the number of participants, their
        interests and prior knowledge.
  \item Only \textbf{Python} may be used for programming.
  \item Implementation, documentation, validation and demonstration examples
        are all part of the project.
  \item Within these rules you are free to choose the implementation, code
        structure and exact content. You can discuss this with me at any time.
\end{itemize}

% ---------------------------------------------------------------------
\section*{Version control and GitHub}
% ---------------------------------------------------------------------
Each group keeps \textbf{one GitHub repository for the whole project}
(Project~2 and the final project), and every member commits from their own
GitHub account.
\begin{itemize}
  \item Use git from the very beginning. Commits should be small, consistent
        and clearly described. Avoid large commits.
  \item Each group member uploads their own code using \textbf{their own
        GitHub account}. Commits from a shared or borrowed account will not
        be credited to the other members.
  \item The repository name is up to your group.
  \item If the repository is private, add my GitHub account
        \textbf{\texttt{say-yas}} as a collaborator.
  \item The repository must contain a README describing the project, how to
        install and run the code, and the AI-use statement (see below).
\end{itemize}

% ---------------------------------------------------------------------
\section*{Use of AI tools}
% ---------------------------------------------------------------------
\textbf{AI-generated code is not accepted in any form.} Code produced by a
chat assistant or code generator (e.g.\ ChatGPT, Claude, Gemini) and copied,
pasted or adapted into your project will be treated as work that is not your
own.
\begin{itemize}
  \item \textbf{Permitted:} inline coding assistance such as GitHub Copilot
        autocompletion while you write the code yourself, \textbf{provided its
        use is acknowledged}.
  \item \textbf{Acknowledgement:} include an AI-use statement in the README
        and in the report that names the tool(s) used and the parts of the
        code (files or functions) where they assisted. If no AI tools were
        used, state this explicitly.
  \item \textbf{Reports and posters must not be generated by AI} in any
        form: their text, figures and layout must be your own work. If AI
        tools were used for any other non-code purpose (e.g.\ checking
        grammar or searching the literature), describe this in the same
        statement.
  \item You are responsible for every line in your repository. In the oral
        exam you must be able to explain and modify any part of the code you
        submitted.
\end{itemize}

% ---------------------------------------------------------------------
\section*{Deliverables and deadlines}
% ---------------------------------------------------------------------
Submit each deliverable following the details provided on the
\textbf{Canvas} page. The written report and poster are handed in on
\textbf{Crowdmark}, and the link to your GitHub repository is submitted
separately on \textbf{Canvas}.

\begin{center}
\small
\renewcommand{\arraystretch}{1.35}
\begin{tabularx}{\textwidth}{|>{\raggedright\arraybackslash}p{2.5cm}|>{\raggedright\arraybackslash}p{3.0cm}|>{\raggedright\arraybackslash}p{3.0cm}|L|}
\hline
\textbf{Deliverable} & \textbf{Due} & \textbf{Submission} & \textbf{Notes} \\
\hline\hline
Group and project title & Monday, October 26, 2026 &
  Group assignment on Canvas &
  Group of two or three members; submit the project title. \\
\hline
Project~2 report & Friday, November 20, 2026 &
  Report (PDF) on Crowdmark; repository link on Canvas &
  3--5 pages; the page limit \textbf{includes} references, figures and
  appendices. \\
\hline
Final project poster & Friday, December 11, 2026 &
  Poster (PDF) on Crowdmark; repository link on Canvas &
  36 inches high by 48 inches wide (landscape). Target audience:
  undergraduate classmates with a basic background in math, physics or
  engineering. \\
\hline
Oral exam & Between November 30 and December 17, 2026, 3:30\,pm &
  Individual meeting; details discussed in class &
  30 minutes per student. Your poster must already be submitted on
  Crowdmark before your exam. The last possible slot ends December 17 at
  3:30\,pm. \\
\hline
\end{tabularx}
\end{center}

\subsection*{Oral exam}
The oral exam is individual and lasts 30 minutes per student. It takes place
between November 30 and December 17, 2026, 3:30\,pm. The oral exams and how
to arrange them will be discussed in class before the Thanksgiving break.

Before your oral exam, your poster must also be submitted on
\textbf{Crowdmark}; that is, you must have finished both Project~2 and the
final project. In the meeting you will:
\begin{itemize}
  \item Present your results and algorithms.
  \item Possibly live-code and run parts of your algorithms to generate
        outputs on the spot.
  \item Answer questions about the theory, your implementation and your
        results. Each member should be able to explain the whole project, not
        only their own parts.
\end{itemize}
Bring a laptop with the repository cloned and a working environment, so that
the code runs without additional setup.

% ---------------------------------------------------------------------
\section*{Grading}
% ---------------------------------------------------------------------
\begin{center}
\small
\renewcommand{\arraystretch}{1.35}
\begin{tabularx}{0.8\textwidth}{|L|>{\raggedleft\arraybackslash}p{3.2cm}|}
\hline
\textbf{Component} & \textbf{Points} \\
\hline\hline
Group and project title (on time) & 10 \\
\hline
Project~2 report & 100 \\
\hline
Final project poster & 100 \\
\hline
Oral exam (individual) & 50 \\
\hline
Bonus: extra steps beyond the required tasks & up to 10 \\
\hline
\end{tabularx}
\end{center}

The Project~2 report and the poster are each graded out of 100. What matters
most is:
\begin{itemize}
  \item \textbf{Coherence:} a clear story from the question, through the
        methods, to the results and conclusions.
  \item \textbf{Correctness:} of the mathematics, the code and the results.
  \item \textbf{Depth of analysis and discussion:} explain what your results
        mean, why they look the way they do, and what their limitations are,
        rather than only showing them.
  \item \textbf{Professional presentation:} figures with labeled axes and
        readable font sizes, captions and legends, a clean and readable
        layout, and proper references.
\end{itemize}

A complete, correct and well-discussed
solution of the required tasks can earn full marks. Tasks marked
\bonus{} and the suggestions under ``Going further'' are optional; exploring
such extra steps can earn \textbf{up to 10 bonus points} in total.

% ---------------------------------------------------------------------
\section*{Requirements for the report and the poster}
% ---------------------------------------------------------------------
Both the report and the poster must contain a clear description of the
project, references, properly labeled figures and an author-contributions
statement.

\subsection*{Content}
\begin{itemize}
  \item A clear and concise description of the problem, the methods and the
        results.
  \item The AI-use statement described above.
  \item The report should not contain any significant amount of code; refer
        to the repository instead.
\end{itemize}

\subsection*{Poster}
\begin{itemize}
  \item The poster must be \textbf{self-contained}: a reader should be able to
        follow it without the report, the code or a spoken explanation.
  \item You may be selective in what you show and how you narrate your story,
        so that an undergraduate audience with a basic background in math,
        physics or engineering can follow it.
  \item You may choose your own title for the poster; be creative.
\end{itemize}

\subsection*{References}
\begin{itemize}
  \item Include a list of references at the end of the report and the poster.
  \item Cite sources clearly where they are used in the text (papers, books,
        software, code repositories and data sets).
  \item Use one consistent citation style throughout.
\end{itemize}

\subsection*{Figures}
\begin{itemize}
  \item Every plot must have labeled $x$ and $y$ axes (and $z$ axis for 3D
        plots or color bars), including units where applicable.
  \item Axis labels, tick labels and legends must have a font size similar to
        the surrounding text, so they remain readable without zooming.
  \item Each figure needs a legend when it shows more than one data series,
        and a caption that states what is shown.
  \item Use log scales where appropriate (e.g.\ for convergence plots) and say
        so in the caption.
\end{itemize}

\subsection*{Author contributions}
Each group must include an author-contributions statement in the report and
the README that states:
\begin{itemize}
  \item Who did what: which parts of the code, experiments, analysis, writing
        and poster each member worked on.
  \item To what extent: the approximate share of each member in each part
        (e.g.\ 50\,\% / 50\,\%).
  \item Who took the lead in which parts.
\end{itemize}
The coding and the report and poster preparation must be shared equally
between the members. The commit history should be consistent with the stated
contributions.

% ---------------------------------------------------------------------
\section*{Codebase requirements}
% ---------------------------------------------------------------------
\begin{itemize}
  \item Well structured and documented, both in the code (docstrings,
        comments) and in the repository (README).
  \item Validated: show that your implementation produces correct results,
        e.g.\ on examples with known solutions.\end{itemize}

% ---------------------------------------------------------------------
\section*{Project description}
% ---------------------------------------------------------------------
\subsection*{Goal}
The goal of this project is to compute mathematical constants and functions
with \emph{continued fractions}. You will implement a general continued
fraction evaluator in Python and use it to compute \textbf{at least 5
constants} and \textbf{at least 5 functions}. You will compare the constants
with reference values from NumPy or SymPy, plot the functions against the
corresponding built-in functions in NumPy, and discuss what you observe: the
benefits and the shortcomings of computing constants and functions this way.

\subsection*{Background}
A generalized continued fraction has the form
\begin{equation*}
  b_0 + \cfrac{a_1}{b_1 + \cfrac{a_2}{b_2 + \cfrac{a_3}{b_3 + \dotsb}}}
\end{equation*}
where the partial numerators $a_n$ and partial denominators $b_n$ are numbers,
or functions of $n$ and of a variable $x$. Truncating after $N$ terms gives the
$N$-th convergent, which can be evaluated from the bottom up (backward
recurrence) or with the three-term recurrences $p_n = b_n p_{n-1} + a_n
p_{n-2}$ and $q_n = b_n q_{n-1} + a_n q_{n-2}$ (forward recurrence).

Continued fractions have a long history \cite{baez}. A classical example is
Brouncker's formula \cite{brouncker},
\begin{equation*}
  \frac{4}{\pi} = 1 + \cfrac{1^2}{2 + \cfrac{3^2}{2 + \cfrac{5^2}{2 + \dotsb}}}.
\end{equation*}
Euler showed that every infinite series can be rewritten as a continued
fraction \cite{euler}, and Gauss found continued fractions for ratios of
hypergeometric functions, which give continued fractions for many elementary
functions \cite{baez}. More recently, the Ramanujan Machine project has used
computer searches to find new \emph{polynomial} continued fractions, where
$a_n$ and $b_n$ are polynomials in $n$, for constants such as $\pi$, $e$,
$\zeta(3)$ and Catalan's constant $G$ \cite{nature,results}. For example, one
conjectured formula is \cite{catalan}
\begin{equation*}
  \frac{1}{2G} = 1 - \cfrac{2}{7 - \cfrac{32}{19 - \cfrac{162}{37 - \dotsb}}},
  \qquad a_n = -2n^4, \quad b_n = 3n^2 + 3n + 1.
\end{equation*}
Many of these formulas are still unproven, so a numerical check is
meaningful in itself.

Different formulas converge at very different speeds: some gain several
digits per term, others need thousands of terms for a few digits. Double
precision (\texttt{float64}) also limits the accuracy to about 16 digits, so
for high-accuracy comparisons you may use exact rational arithmetic (Python's
\texttt{fractions} module \cite{fractions}) or arbitrary precision
(\texttt{mpmath} \cite{mpmath}, which SymPy uses internally).

\subsection*{Tasks for Project 2 (report, due November 20, 2026)}
\begin{enumerate}
  \item Implement a general continued fraction evaluator that takes $a_n$ and
        $b_n$ as functions of $n$ and a truncation depth $N$, using the
        backward recurrence.
  \item Implement \textbf{at least 5 constants} with continued fraction
        formulas, for example $\pi$, $e$, $\sqrt{2}$, the golden ratio,
        $\ln 2$, $\zeta(2)$, $\zeta(3)$ and Catalan's constant. Use the
        sources below \cite{brouncker,baez,catalan,euler,results} and cite the
        source of every formula you use.
  \item Implement \textbf{at least 5 functions} with continued fraction
        formulas, for example $e^x$, $\ln(1+x)$, $\arctan x$, $\tan x$,
        $\tanh x$ and $\sqrt{1+x}$, and cite your sources.
  \item Compare each constant with a reference value from NumPy or SymPy,
        and report the absolute and relative error in a table.
  \item Plot each function against the corresponding built-in function in
        NumPy over a suitable interval.
  \item \bonus\ Also implement the forward recurrence and check that it
        agrees with the backward recurrence.
\end{enumerate}

\subsection*{Tasks for the final project (poster, due December 11, 2026)}
\begin{enumerate}
  \item Compare and plot the accuracy of your results: for the constants,
        plot the error against the depth $N$; for the functions, plot the
        error against $x$ for a few depths $N$ (log scales).
  \item Discuss what you observe: how fast does each formula converge, how
        does the accuracy of the functions depend on $N$ and on $x$, and
        where does the continued fraction converge well and where not? What
        are the benefits and the shortcomings of computing constants and
        functions with continued fractions?
  \item \bonus\ Compare how fast different formulas for the same constant
        converge.
  \item \bonus\ Compare with a truncated Taylor series of the same
        function.
\end{enumerate}

\subsection*{Going further}
These optional steps count toward the bonus points. You are welcome to go
beyond the provided links, for example by using Euler's formula \cite{euler}
to turn known series into new continued fractions, or by computing best
rational approximations such as $355/113$ for $\pi$.

% ---------------------------------------------------------------------
\begin{thebibliography}{9}
  \bibitem{brouncker} Wikipedia. \emph{William Brouncker, 2nd Viscount
    Brouncker}, section ``Brouncker's formula''.
    \url{https://en.wikipedia.org/wiki/William_Brouncker,_2nd_Viscount_Brouncker}
  \bibitem{baez} J.~Baez (2020). \emph{Three Phases of Continued Fraction
    Theory}. The $n$-Category Caf\'e, September 6, 2020.
    \url{https://golem.ph.utexas.edu/category/2020/09/five_levels_of_continued_fract.html}
  \bibitem{catalan} The Ramanujan Machine (2020). \emph{Catalan's constant
    $G$ results}.
    \url{https://ramanujanmachine.com/wp-content/uploads/2020/06/catalan.pdf}
  \bibitem{euler} The Ramanujan Machine. \emph{From polynomial continued
    fraction to infinite sums}.
    \url{https://ramanujanmachine.com/the-mathematics-of-polynomial-continued-fractions/from-polynomial-continued-fraction-to-infinite-sums/}
  \bibitem{results} The Ramanujan Machine. \emph{Results} (with links to
    results for $\pi$, $e$, $\ln 2$, $\zeta(2)$, $\zeta(3)$ and Catalan's
    constant). \url{https://ramanujanmachine.com/results/}
  \bibitem{nature} G.~Raayoni, S.~Gottlieb, Y.~Manor, G.~Pisha, Y.~Harris,
    U.~Mendlovic, D.~Haviv, Y.~Hadad and I.~Kaminer (2021). Generating
    conjectures on fundamental constants with the Ramanujan Machine.
    \emph{Nature}, 590, 67--73. \url{https://doi.org/10.1038/s41586-021-03229-4}
  \bibitem{fractions} Python Software Foundation. \emph{fractions ---
    Rational numbers}, Python documentation.
    \url{https://docs.python.org/3/library/fractions.html}
  \bibitem{mpmath} \emph{mpmath: a Python library for arbitrary-precision
    floating-point arithmetic}. \url{https://mpmath.org}
\end{thebibliography}

\end{document}
